% AUTO-GENERATED by paper/naacl2027/make_tables.py -- do not hand-edit.
% Source: paper/results/corpus_page_ranges.json
% Regenerate with: python paper/naacl2027/make_tables.py
\begin{table*}[t]
\centering
\small
\begin{tabular}{llrp{0.62\textwidth}}
\toprule
\textbf{State} & \textbf{Tier} & \textbf{pp} & \textbf{Retained PDF pages} \\
\midrule
AZ & trimmed & 65 & 2, 21, 23--28, 41, 43--50, 66, 70--82, 107, 110--120, 139, 142--145, 156, 158--162, 179, 181--185, 198, 201--204 (of 217) \\
AZ & subset & 15 & 2, 21, 23, 66, 70--78, 81--82 (of 217) \\
CA & subset & 13 & 1--5, 11--13, 18--22 (of 68) \\
CO & trimmed & 41 & 1--41 (of 41) \\
CO & subset & 10 & 1, 3--9, 11--12 (of 41) \\
TX & trimmed & 71 & 1, 9--65, 69--81 (of 87) \\
TX & subset & 9 & 1, 9--13, 59--61 (of 87) \\
NV & trimmed & 71 & 1, 3, 16--84 (of 98) \\
NV & subset & 15 & 1, 3, 23--26, 54--57, 61--65 (of 98) \\
KY & trimmed & 52 & 1, 52--102 (of 120) \\
KY & subset & 8 & 1, 52--55, 65, 71--72 (of 120) \\
\bottomrule
\end{tabular}
\caption{Pages each tier retains, as indices into the published PDF file. Printed page numbers differ by each document's front matter (Kentucky's PDF page 52 is printed page 39). Every tier page was matched by text to a published page. California has no trimmed tier.}
\label{tab:corpus-pages}
\end{table*}
